% =====================================================================
% 1. Introdução
% =====================================================================
\section{Introdução}\label{sec:introducao}

O pipeline gráfico é a sequência de etapas que transforma objetos 3D em uma imagem 2D: modelagem,
composição da cena, câmera, projeção, iluminação e rasterização \cite{azevedo2018}.
\confirmar{edição, ano e editora do livro Azevedo; Conci; Vasconcelos} Neste trabalho, o grupo
percorreu todas as etapas com cinco objetos do cotidiano: lápis, peão de xadrez, haltere, martelo e
lixeira. Cada objeto foi descrito por uma função matemática e convertido em malha de triângulos com
o algoritmo \textit{Marching Cubes} \cite{lorensen1987}. A cena foi vista por uma câmera,
projetada em perspectiva, iluminada pelo modelo de Phong \cite{phong1975} e rasterizada por um
rasterizador implementado pelo grupo.

O código foi escrito em Python. A biblioteca NumPy \cite{harris2020} cuidou das operações com
matrizes, e a Matplotlib \cite{hunter2007} exibiu as imagens. Apenas o \textit{Marching Cubes}
veio pronto, da biblioteca scikit-image \cite{vanderwalt2014}. As demais etapas foram
implementadas pelo grupo.
